\documentclass[tikz,border=15pt]{standalone}
\usepackage{fontspec}
\usepackage{polyglossia}
\setmainlanguage{russian}
\setotherlanguage{english}

% Системный шрифт Times New Roman (Windows)
\setmainfont{Times New Roman}
\setsansfont{Arial}

\usepackage{amsmath}
\usetikzlibrary{arrows.meta,positioning,calc}

\begin{document}
\begin{tikzpicture}[
    x=1.4cm, y=0.0028cm,   % масштаб: 1 ед. по X = 1.4 см, по Y = 0.0028 см/кэВ
    >=Stealth,
    every node/.style={font=\small},
    beta/.style={->, line width=1.6pt, red!70!black},
    gamma/.style={->, line width=2.5pt, blue!70!black},
    level/.style={line width=1.4pt},
    lbl/.style={font=\footnotesize},
    axis/.style={->, line width=1.0pt}
]

% ==================== ОСИ ====================
% Ось Y (энергия) слева
\draw[axis] (0,0) -- (0,3700) node[above] {\(E\), кэВ};
% Ось X (условная позиция)
\draw[axis] (0,0) -- (12,0);

% Тики на оси Y (слева, внутрь)
\foreach \e in {0, 500, 1000, 1500, 2000, 2500, 3000, 3500} {
    \draw[axis] (0,\e) -- (0.15,\e) node[left] {\e};
}

% ==================== УРОВНИ 82Kr ====================
% Основное состояние
\draw[level] (5,0) -- (11,0) node[right] {0 \quad (\(0^+\))};
\node[below] at (8,-100) {\(^{82}\)Kr (стабильный)};

% Возбуждённые состояния
\draw[level] (5,776.5)  -- (11,776.5)  node[right] {776.5 \quad (\(2^+\))};
\draw[level] (5,1474.8) -- (11,1474.8) node[right] {1474.8 \quad (\(2^+\))};
\draw[level] (5,2094.0) -- (11,2094.0) node[right] {2094.0 \quad (\(3^+\))};
\draw[level] (5,3181.1) -- (11,3181.1) node[right] {3181.1 \quad (\(2^+\))};

% ==================== РОДИТЕЛЬСКИЙ УРОВЕНЬ 82Br ====================
\draw[level] (1.5,2094.0) -- (3.5,2094.0) node[midway,above] {\(^{82}\)Br (\(2^-\))};
\node[lbl, below] at (2.5,2000) {\(T_{1/2}=35.28\) ч};

% ==================== БЕТА-ПЕРЕХОДЫ ====================
% Ветвь на 3181.1 кэВ (E_max = 444 кэВ, ~80%)
\draw[beta] (3.5,2094) -- (4.8,3181.1) node[midway,above,sloped] {\(\beta^-\)};
\node[lbl, right, align=left] at (3.6,2900) {\(E_{\max}=444\) кэВ \\ \(\sim80\%\)};

% Ветвь на 1474.8 кэВ (E_max = 1535 кэВ, ~18%)
\draw[beta] (3.5,2094) -- (4.8,1474.8) node[midway,above,sloped] {\(\beta^-\)};
\node[lbl, right, align=left] at (3.6,1750) {\(E_{\max}=1535\) кэВ \\ \(\sim18\%\)};

% ==================== ГАММА-ПЕРЕХОДЫ ====================
% 1317.5 кэВ (2094.0 -> 776.5)
\draw[gamma] (6.0,2094.0) -- (6.0,776.5);
\node[text=blue!70!black, fill=white, inner sep=1pt, rotate=90] at (5.85,1435) {1317.5};

% 776.5 кэВ (776.5 -> 0)
\draw[gamma] (7.5,776.5) -- (7.5,0);
\node[text=blue!70!black, fill=white, inner sep=1pt, rotate=90] at (7.35,388) {776.5};

% 1474.8 кэВ (1474.8 -> 0)
\draw[gamma] (9.5,1474.8) -- (9.5,0);
\node[text=blue!70!black, fill=white, inner sep=1pt, rotate=90] at (9.35,737) {1474.8};

% 1044.0 кэВ (2094.0 -> 776.5, упрощённо)
\draw[gamma] (8.5,2094.0) -- (8.5,776.5);
\node[text=blue!70!black, fill=white, inner sep=1pt, rotate=90] at (8.35,1435) {1044.0};

% 619.1 кэВ (2094.0 -> 1474.8)
\draw[gamma] (5.3,2094.0) -- (5.3,1474.8);
\node[text=blue!70!black, fill=white, inner sep=1pt, rotate=90] at (5.15,1784) {619.1};

% ==================== ЗАГОЛОВКИ ====================
\node[anchor=south east, font=\footnotesize\bfseries] at (11,3700) {\(J^\pi\), \(E\), кэВ};
\node[font=\large\bfseries, anchor=west] at (0,3700) {Упрощённая схема распада \(^{82}\)Br};

\end{tikzpicture}
\end{document}